Broadly, the three cases tell different stories.

Qwen finetunes
The finetuned Qwen models are very close to base in hidden space: cosine stays around 0.997-0.998 through most layers, with larger L2 only toward the last layer. So the representation drift is subtle, not a wholesale rewrite.

The logit-side differences are also mostly small through early/mid layers, then become more visible at the last layer. Raw JSD/TVD and Jaccard degradation are strongest at Last, which suggests the behavioral difference is concentrated near final readout rather than distributed evenly through the whole stack. Model-norm makes the layer trend smoother and shows accuracy rising late, as expected.

Quantized Llama
This is the clearest “severity ladder”:

8-bit is close to base.
4-bit is further but still structurally similar.
HF1Bit is qualitatively much more distorted.

Hidden cosine and L2 show that directly: 8-bit remains high-cosine / low-L2, 4-bit worsens, HF1Bit has much lower cosine and much larger L2. Logit metrics follow the same pattern: TVD/JSD are small for 8-bit, larger for 4-bit, and huge for HF1Bit. Jaccard@k collapses especially for HF1Bit.

Layerwise, quantization errors tend to grow toward later layers, especially in raw lens. That means small perturbations early become more behaviorally consequential closer to output.

Pythia training dynamics
Pythia shows the strongest developmental signal.

Early checkpoints, 1k, are far from the final step143000 model across hidden and logit metrics. Mid checkpoints, 71k, are consistently closer to final than 1k, especially in larger models.

The larger Pythia models show a more interpretable depth progression: early/mid layers can already partially align, but late layers and last-layer behavior distinguish training stage strongly. In the mid checkpoint, larger models like 6.9B and 12B have much lower TVD/JSD and higher Jaccard than their 1k versions, suggesting they have moved much closer to the final model’s next-token geometry.

The base/comp accuracy rows also make sense: final/reference-side accuracy rises strongly late; early checkpoints have weaker late-layer accuracy, while 71k checkpoints look much closer.

Correlation Tables
The correlations mostly answer: “which intermediate layers already predict the last-layer behavior?”

High correlations near Layer/Last close to the final layer mean the metric has stabilized late in the network. For smooth metrics like hidden L2/cosine, strongest correlations are often near the last layer. For discrete metrics like Jaccard or accuracy, correlations can be noisier or undefined when values are constant.

So the overall picture is:

Qwen finetuning: small, late-concentrated behavioral drift.
Llama quantization: monotonic degradation by quantization severity, strongest for HF1Bit.
Pythia training: early checkpoints are genuinely different systems; mid checkpoints become increasingly final-like, especially in larger models and later layers.




9:20 PM
and also for the raw and mdoel norm lens
9:22 PM


Yes. The raw vs model-norm lens difference is important.

Raw Lens
Raw uses the hidden state directly through the unembedding/head. It answers: “If I project this intermediate residual state as-is, how similar are the two models?”

So raw is more sensitive to scale, norm drift, and representational calibration differences across layers/models. In the tables, raw often shows:

larger TVD/JSD,
lower Jaccard,
worse early-layer accuracy,
stronger late-layer jumps.
For Qwen, raw differences stay small early/mid but spike at Last, especially JSD/TVD. That suggests the finetune does not radically change most intermediate representations, but the final prediction distribution still shifts.

For Llama quantization, raw clearly separates 8-bit, 4-bit, and HF1Bit. HF1Bit raw TVD/JSD and Jaccard are very poor, which means the unnormalized hidden states are not projecting into base-like logits.

For Pythia, raw is the more severe developmental comparison: 1k checkpoints are far from step143000, and 71k checkpoints are closer but still size/layer dependent.

Model-Norm Lens
Model-norm applies the saved final normalization before projection. It asks: “If this intermediate hidden state were normalized like the final layer, how similar are the logits?”

This often reduces pure scale/norm artifacts and makes the comparison more about directional/content alignment. In practice:

Qwen model-norm has smoother TVD/JSD than raw.
Accuracy becomes much more meaningful across depth, rising toward late/last layers.
Jaccard is often more stable in model-norm than raw.
For Llama, model-norm can expose whether quantization preserved semantic direction even if raw scale is off.
For Pythia, model-norm helps distinguish “the representation is wrong” from “the representation is not yet calibrated for the head.”
Interpretation
If raw differs a lot but model-norm is closer, the models may share similar representational directions but differ in activation scale/calibration.

If both raw and model-norm differ a lot, the actual representational content is different.

So:

Qwen: mostly shared representation; differences are late and partly readout/calibration-sensitive.
8-bit Llama: close under both, so quantization preserves content and calibration well.
4-bit Llama: moderate degradation under both.
HF1Bit: bad under both, so not just a norm/calibration issue.
Pythia 1k: far under both, so early checkpoints are genuinely unlike the final model.
Pythia 71k: closer, especially in larger models, but late-layer behavior still shows where training has not fully converged.